\documentclass[11pt,t,aspectratio=169]{beamer}
\setbeamersize{text margin left=1em,text margin right=1em}
\usetheme{sussexdark} % options: [footlogo=false] hides the footer logo
\usepackage{hyperref}
\usepackage{physics}
\usepackage{multicol}
\usepackage{listings}

% code listings in Sussex colours
\lstset{
    basicstyle=\ttfamily,
    keywordstyle=\color{usalbionblue}\bfseries,
    commentstyle=\color{usditchlinggreen}\itshape,
    stringstyle=\color{uspavilionpurple},
    columns=fullflexible,
    keepspaces=true,
    showstringspaces=false,
    aboveskip=2pt,
    belowskip=0pt,
}

% Delete this, if you do not want the table of contents to pop up at the beginning of each subsection:
\AtBeginSection[]
{
    \begingroup
        \setbeamertemplate{background canvas}[vertical shading][bottom=usgreen,top=usgreen]
        \setbeamertemplate{footline}[sectionfootline]
        \setbeamertemplate{section page}[mysection]
        \frame[c]{
        \sectionpage
        }
    \endgroup
}

\title{16:9 Dark Beamer Template for the University of Sussex}
\subtitle{Not an official University of Sussex template}
\author{Debby Oluwatuyi}
\institute{University of Sussex}
\date{\today}

%titlepage logo
\titlegraphic{\includegraphics[height=0.95in]{logos/sussex-logo-chalk.pdf}}

% colour swatch for the palette slide
\newcommand{\swatch}[3]{%
    \begin{minipage}[t]{0.118\textwidth}\centering
        {\color{#1}\rule{\linewidth}{0.3cm}}\par\vspace{1pt}
        {\tiny\sussexsemibold #2\par}
        {\fontsize{4.5}{5}\selectfont\texttt{#1}\par \#{}#3\par}
    \end{minipage}\hfill}

\begin{document}

\begin{frame}
    \titlepage
\end{frame}

\begin{frame}[c]{Outline}
    \tableofcontents
\end{frame}

\section{Disclaimer}
\begin{frame}[c]
\frametitle{Disclaimer}
    \begin{multicols}{2}
    \begin{itemize}
        \item THIS IS NOT AN OFFICIAL UNIVERSITY OF SUSSEX TEMPLATE.
        \item Adapted from the \LaTeX{} template by Guanyang Xue (itself modified from Alex Pacheco's).
        \item Follows the University of Sussex Brand Guidelines V.1 (October 2025).
        \item Modify if you have suggestions to improve it and pass on the modifications so that anyone interested in using \LaTeX{} for presentations could benefit.
        \end{itemize}
    \end{multicols}
    \begin{itemize}
        \item Typing equations with \emph{physics} package is always more convenient.
        $$\pdv{\vb*{U}}{t}+(\vb*{U}\cdot\nabla)\vb*{U}-\nu\laplacian\vb*{U} =-\frac{1}{\rho }\grad p+\vb*{g}$$
    \end{itemize}
\end{frame}

\section{How to}
\begin{frame}[c,fragile]
    \frametitle{How to use it}
    \framesubtitle{How to install it to your path}
    \begin{itemize}
        \item Place the contents within your TEXINPUTS path. On Linux, modify the TEXINPUTS variable to point to the directory where these files exist.
        \item For example: \lstinline[language=bash]|export TEXINPUTS=${HOME}/LaTeX//:${TEXINPUTS}|
        \item On MacTeX, you need to install it to \lstinline|${HOME}/Library/texmf/tex/latex/sussexdark|
        \item If you have admin permission copy it to your \TeX{} tree to the directory where the other beamer themes are located e.g. \lstinline|/usr/share/texmf-texlive/tex/latex/beamer/themes|
        \item To use the theme, add the command \lstinline[language={[LaTeX]TeX}]|\usetheme{sussexdark}| to the preamble of your document.
        \item Compile using pdflatex to get the pdf file.
        \item Feel free to modify the theme to suit your needs.
    \end{itemize}
\end{frame}
\subsection{Blocks}
\begin{frame}[c,fragile]
    \frametitle{How to use it}
    \framesubtitle{How to use blocks}
    \begin{block}{This is a block}
\begin{lstlisting}[language={[LaTeX]TeX}]
\usepackage{listings}
...
\begin{frame}[c,fragile] % Use fragile option
\end{lstlisting}
    \end{block}
    \begin{exampleblock}{This is an example block}
\begin{lstlisting}[language=Python]
import numpy as np
a = np.array([1, 2, 3])
\end{lstlisting}
    \end{exampleblock}
    \begin{alertblock}{This is an alert block}
\begin{lstlisting}[language=bash]
fatal error: mpi.h: No such file or directory
compilation terminated.
\end{lstlisting}
    \end{alertblock}
\end{frame}

\subsection{Colours}
\begin{frame}[c]
    \frametitle{Brand colours}
    \framesubtitle{Use them as \texttt{\textbackslash color\{usgreen\}}; official tints exist as \texttt{usgreen80}, \texttt{uschalk40}, \ldots}
    {\footnotesize\sussexsemibold Sussex by the Sea (primary)}\par\smallskip
    \swatch{usgreen}{South Downs Green}{033803}
    \swatch{uschalk}{Seven Sisters Chalk}{EAE5DD}
    \swatch{usturquoise}{English Channel Turquoise}{32D8C5}
    \hspace*{0.6\textwidth}\par\smallskip
    {\footnotesize\sussexsemibold Brighton (secondary)}\par\smallskip
    \swatch{usargusred}{Argus Red}{FF3323}
    \swatch{usnorthlaineorange}{North Laine Orange}{FF9200}
    \swatch{uscarouselyellow}{Carousel Yellow}{FEFE00}
    \swatch{usrockpink}{Rock Pink}{FF00C5}
    \swatch{uspavilionpurple}{Pavilion Purple}{CA21E8}
    \swatch{usalbionblue}{Albion Blue}{0713CB}
    \swatch{usbeachhutblue}{Beach Hut Blue}{00FFFF}
    \swatch{usstanmergreen}{Stanmer Green}{0FED8D}\par\smallskip
    {\footnotesize\sussexsemibold South Downs (secondary)}\par\smallskip
    \swatch{usshellpeach}{Shell Peach}{F5C48A}
    \swatch{usheatherpurple}{Heather Purple}{5B3046}
    \swatch{usrampionlilac}{Rampion Lilac}{A193F2}
    \swatch{usflintgrey}{Flint Grey}{6F7587}
    \swatch{usoceanteal}{Ocean Teal}{003B4A}
    \swatch{usadonisblue}{Adonis Blue}{3F77CA}
    \swatch{usditchlinggreen}{Ditchling Green}{6B8400}
    \swatch{usdownlandgreen}{Downland Green}{A7D936}
\end{frame}
\end{document}
